\chapter{Arsitektur Bisnis Target untuk Transformasi Digital}

\section*{Tujuan Pembelajaran}
\begin{itemize}
  \item Menganalisis pilihan transformasi digital memakai kerangka \textit{Designed for Digital}.
  \item Merumuskan model operasi target dan kapabilitas yang perlu dibangun, diperkuat, atau dihentikan.
  \item Merancang \textit{value stream} target beserta ukuran keberhasilannya.
\end{itemize}

\section{Pendahuluan}

% TODO
% Cakupan: transformasi digital yang dimulai dari membeli teknologi tanpa mengubah cara kerja,
%   sehingga sistem baru hanya memindahkan masalah lama ke layar baru.

\section{Transformasi Digital dari Sudut Pandang Arsitektur}

% TODO
% Cakupan: beda digitisasi, digitalisasi, dan transformasi digital; lima blok bangunan Designed for
%   Digital (shared customer insights, operational backbone, digital platform, accountability
%   framework, external developer platform); beda operational backbone dan digital platform.
% Kasus Medikara: Medikara+ sebagai calon platform digital dan kebutuhan backbone yang rapi lebih
%   dahulu.

\section{Model Operasi Target}

% TODO
% Cakupan: pergeseran model operasi dari kondisi saat ini; standardisasi versus keleluasaan lokal;
%   dampak pada struktur organisasi, data bersama, dan aplikasi bersama.

\section{Opsi Kondisi Target dan Business Scenario}

% TODO
% Cakupan: menyusun dua sampai tiga opsi; kriteria evaluasi yang diturunkan dari prinsip C1;
%   business scenario dalam fase B; mendokumentasikan alasan memilih dan hal yang dikorbankan.
% Kasus Medikara: opsi sentralisasi penuh layanan pasien versus inti yang distandarkan dengan
%   keleluasaan lokal.

\section{Kapabilitas Target dan Peta Perubahan Kapabilitas}

% TODO
% Cakupan: kapabilitas baru, diperkuat, dan dihentikan; heatmap target; kapabilitas pengelolaan
%   ekosistem mitra (model bisnis platform dibahas rinci di mata kuliah IT40613).
% Pendalaman teknis: memilih kapabilitas yang layak menjadi layanan bersama dan menuliskan antarmuka
%   bisnisnya sebagai bahan desain API.

\section{Value Stream Target}

% TODO
% Cakupan: merancang ulang tahapan dari sudut pandang pasien; ukuran keberhasilan per tahap;
%   keterkaitan tahap dengan kapabilitas target.
% Kasus Medikara: value stream rawat jalan dengan satu identitas dan satu rekam medis lintas
%   fasilitas.

\section{Aktivitas: Memilih Opsi Kondisi Target Bisnis}

% TODO
% Tujuan aktivitas: memilih arah arsitektur bisnis target dengan alasan yang dapat dipertanggungjawabkan.
% Peran dalam tim: analis bisnis memimpin perumusan opsi dan kriteria; analis teknologi menilai
%   kelayakan tiap opsi dari kondisi saat ini; pemodel memperbarui capability map; koordinator
%   menuliskan keputusan dan alasannya.
% Langkah: susun dua sampai tiga opsi; tetapkan kriteria berbobot; nilai tiap opsi; pilih satu dan
%   tulis trade-off-nya; tandai perubahan kapabilitas; rancang value stream target.
% Luaran dan tonggak capstone: model operasi target, capability map target, dan value stream
%   target (C3).
% Berkas sources yang direncanakan: sources/chapter08/templat-opsi-kondisi-target.md,
%   sources/chapter08/templat-matriks-kriteria-opsi.csv,
%   sources/chapter08/medikara-capability-map-target.csv

\section{Ringkasan}

% TODO

\section{Latihan}

\begin{enumerate}
  \item Jelaskan perbedaan \textit{operational backbone} dan \textit{digital platform} dengan contoh dari Medikara.
  \item Sentralisasi layanan pasien memberi efisiensi tetapi mengurangi keleluasaan rumah sakit di daerah. Tuliskan \textit{trade-off} ini dalam tiga kalimat yang dapat dibaca direksi.
  \item Mengapa kapabilitas pengelolaan kemitraan muncul di kondisi target tetapi tidak di kondisi saat ini? Apa syarat agar kapabilitas itu berjalan?
\end{enumerate}
